\section*{Reading list by section}
The papers behind each part of the deck, in the order they are discussed, with one line on why each is worth reading. Entries with an arXiv identifier are included in the paper archive; entries without one are listed with a URL. Items marked \emph{secondary source} are blog posts, model cards or commentary.

\subsection*{1. Overview and the long view (slides 1-4)}
\begin{itemize}[leftmargin=1.2em,itemsep=3pt]\sloppy
\item \textbf{Dally, Keckler and Kirk (2021).} \textit{Evolution of the graphics processing unit (GPU)}. \url{https://doi.org/10.1109/MM.2021.3113475}. IEEE Micro 41(6). The architects' own twenty-year retrospective, from G80 to Ampere: why unified shaders, why SIMT, why the datapath specialized. The best single companion to this deck.
\item \textbf{Nickolls and Dally (2010).} \textit{The GPU computing era}. \url{https://doi.org/10.1109/MM.2010.41}. IEEE Micro 30(2). The state of the art at Fermi's launch: the Tesla line from G80 to GF100, the programming model, and the case that GPUs had become general processors.
\item \textbf{Garland and Kirk (2010).} \textit{Understanding throughput-oriented architectures}. \url{https://doi.org/10.1145/1839676.1839694}. CACM 53(11), open access. The clearest statement of the latency-versus-throughput framing that slide 2 uses: many threads, simple cores, wide memory.
\item \textbf{Fatahalian and Houston (2008).} \textit{A closer look at GPUs}. \url{https://doi.org/10.1145/1400181.1400197}. CACM 51(10). A short, precise account of how a GPU core differs from a CPU core (SIMD lanes, multithreading, shared state), written as G80 and CUDA arrived.
\item \textbf{Owens, Houston, Luebke, Green, Stone and Phillips (2008).} \textit{GPU computing}. \url{https://doi.org/10.1109/JPROC.2008.917757}. Proceedings of the IEEE 96(5); open copy at escholarship.org/content/qt0cv1p1nc/qt0cv1p1nc.pdf. The survey written at the moment GPGPU became GPU computing: hardware, programming models, applications.
\item \textbf{Luebke and Humphreys (2007).} \textit{How GPUs work}. \url{https://doi.org/10.1109/MC.2007.59}. IEEE Computer 40(2). A four-page primer on the pipeline and its programmable stages, for readers who need the graphics side of slides 5-6 explained from zero.
\item \textbf{Peddie (2022).} \textit{The History of the GPU: Steps to Invention}. \url{https://doi.org/10.1007/978-3-031-10968-3}. Springer, volume 1 of three. The pre-history: from early raster hardware to the first consumer 3D chips; dense with dates and die photographs.
\item \textbf{Peddie (2022).} \textit{The History of the GPU: Eras and Environment}. \url{https://doi.org/10.1007/978-3-031-13581-1}. Springer, volume 2. 3dfx to unified shaders; the industrial history behind the timeline on slide 4.
\item \textbf{Peddie (2023).} \textit{The History of the GPU: New Developments}. \url{https://doi.org/10.1007/978-3-031-14047-1}. Springer, volume 3. The 2010s and the turn to AI; overlaps the last third of this deck.
\item \textbf{Hennessy and Patterson (2017).} \textit{Computer Architecture: A Quantitative Approach, 6th edition, chapter 4 (data-level parallelism in vector, SIMD and GPU architectures)}. \url{https://www.elsevier.com/books/computer-architecture/hennessy/978-0-12-811905-1}. Textbook. Chapter 4 unifies vector machines, SIMD extensions and GPUs in one vocabulary and works a Pascal example; the 5th edition (2011) does the same with Fermi.
\end{itemize}
\subsection*{2. The graphics pipeline, fixed and programmable (slides 5-6)}
\begin{itemize}[leftmargin=1.2em,itemsep=3pt]\sloppy
\item \textbf{Clark (1982).} \textit{The Geometry Engine: a VLSI geometry system for graphics}. \url{https://doi.org/10.1145/800064.801272}. SIGGRAPH 1982; open copy at web.stanford.edu/class/cs349g/p127-clark.pdf. The transform stage of the pipeline as a custom chip: the origin of fixed-function graphics hardware and of Silicon Graphics.
\item \textbf{Akeley (1993).} \textit{RealityEngine graphics}. \url{https://doi.org/10.1145/166117.166131}. SIGGRAPH 1993. SGI's high-end pipeline built from boards of Geometry Engines and raster chips; the fully fixed-function pipeline at its peak.
\item \textbf{Smith (The Register) (1999).} \textit{Nvidia unveils GeForce 256 (August 31, 1999)}. \url{https://www.theregister.com/1999/08/31/nvidia_unveils_256cpu_cray_geforce/}. Press, secondary source. The launch coverage with NVIDIA's claims as made: 'nearly 23 million transistors', 480 Mpixels/s, and the coining of 'graphics processing unit'. Component databases list 17 M transistors.
\item \textbf{Lindholm, Kilgard and Moreton (2001).} \textit{A user-programmable vertex engine}. \url{https://doi.org/10.1145/383259.383274}. SIGGRAPH 2001. The GeForce 3's programmable vertex shader: the instruction set, the register model, and why the stage was made programmable. Slide 6.
\item \textbf{Mark, Glanville, Akeley and Kilgard (2003).} \textit{Cg: a system for programming graphics hardware in a C-like language}. \url{https://doi.org/10.1145/882262.882362}. ACM TOG 22(3), SIGGRAPH 2003. The shading language that made shaders look like C and whose authors went on to CUDA; read with Kilgard's two-page summary, arXiv cs/0302013.
\item \textbf{Kilgard (2003).} \textit{Cg in two pages}. arXiv:cs/0302013. A compact statement of the Cg programming model; the arXiv copy is the easiest to obtain.
\item \textbf{Montrym and Moreton (2005).} \textit{The GeForce 6800}. \url{https://doi.org/10.1109/MM.2005.37}. IEEE Micro 25(2). The last large pre-unified NVIDIA design: separate vertex and pixel processors, Shader Model 3.0, and the load-balancing problem that G80 solved.
\item \textbf{Andrews and Baker (2006).} \textit{Xbox 360 system architecture}. \url{https://doi.org/10.1109/MM.2006.45}. IEEE Micro 26(2). The Xenos GPU (ATI, 2005), the first unified-shader design to ship, a year before G80.
\item \textbf{Blythe (2006).} \textit{The Direct3D 10 system}. \url{https://doi.org/10.1145/1141911.1141947}. SIGGRAPH 2006. The API side of unification: Shader Model 4.0's common shader core and why the fixed distinction between vertex and pixel programs was dropped.
\item \textbf{Peddie (IEEE Computer Society, Chasing Pixels) (2021).} \textit{AMD's unified shader GPU history}. \url{https://www.computer.org/publications/tech-news/chasing-pixels/amd-unified-shader-gpu-history}. Secondary source. Settles the 'first unified shader' question: Xenos first in a console (2005), G80 first on the PC with a common API (DirectX 10, 2006).
\end{itemize}
\subsection*{3. GPGPU and the unified G80 with CUDA (slides 7-8)}
\begin{itemize}[leftmargin=1.2em,itemsep=3pt]\sloppy
\item \textbf{Buck, Foley, Horn, Sugerman, Fatahalian, Houston and Hanrahan (2004).} \textit{Brook for GPUs: stream computing on graphics hardware}. \url{https://doi.org/10.1145/1015706.1015800}. ACM TOG 23(3), SIGGRAPH 2004; project page graphics.stanford.edu/papers/brookgpu. Streams and kernels compiled to shaders: the programming model that CUDA's first author carried to NVIDIA. Slide 7.
\item \textbf{Owens, Luebke, Govindaraju, Harris, Krueger, Lefohn and Purcell (2007).} \textit{A survey of general-purpose computation on graphics hardware}. \url{https://doi.org/10.1111/j.1467-8659.2007.01012.x}. Computer Graphics Forum 26(1); open copy at escholarship.org/content/qt9ns2d70c/qt9ns2d70c.pdf. The catalog of what ran on shaders before CUDA, and of the tricks (scatter via vertex processors, multi-pass reductions) that CUDA made unnecessary.
\item \textbf{Steinkraus, Simard and Buck (2005).} \textit{Using GPUs for machine learning algorithms}. \url{https://doi.org/10.1109/ICDAR.2005.251}. ICDAR 2005. The earliest paper training neural networks on a GPU, through the graphics API; three of the authors' ideas reappear in cuDNN a decade later.
\item \textbf{Chellapilla, Puri and Simard (2006).} \textit{High performance convolutional neural networks for document processing}. \url{https://inria.hal.science/inria-00112631}. IWFHR 2006, open access. Convolution unrolled into matrix multiply and run on a GPU: the formulation every later framework used.
\item \textbf{Lindholm, Nickolls, Oberman and Montrym (2008).} \textit{NVIDIA Tesla: a unified graphics and computing architecture}. \url{https://doi.org/10.1109/MM.2008.31}. IEEE Micro 28(2). The canonical G80 paper: streaming processors, the SM, SIMT (the term is introduced here), the memory system, and the graphics and compute pipelines sharing one array. Slides 8 and 11.
\item \textbf{Nickolls, Buck, Garland and Skadron (2008).} \textit{Scalable parallel programming with CUDA}. \url{https://doi.org/10.1145/1365490.1365500}. ACM Queue 6(2), open access. The programming model as its designers explain it: threads, blocks, grids, shared memory, barriers, and why blocks are independent. Slide 9.
\item \textbf{Kirk (2007).} \textit{NVIDIA CUDA software and GPU parallel computing architecture (ISMM 2007 keynote)}. \url{https://doi.org/10.1145/1296907.1296909}. The first appearance of CUDA's architecture in the academic literature, four months after the public beta.
\item \textbf{eWeek (2006).} \textit{Nvidia launches new architecture for its graphics chips (November 8, 2006)}. \url{https://www.eweek.com/networking/nvidia-launches-new-architecture-for-its-graphics-chips/}. Press, secondary source. Confirms that the GeForce 8800 and CUDA were announced on the same day.
\item \textbf{NVIDIA (2007).} \textit{Tesla product overview: C870, D870, S870 (June 2007)}. \url{https://www.nvidia.com/docs/io/43395/tesla_product_overview.pdf}. Vendor datasheet. The first compute-only products: 128 thread processors, 1.5 GB, 76.8 GB/s, 'up to 518 gigaflops'.
\item \textbf{Khronos Group (2008).} \textit{The Khronos Group releases OpenCL 1.0 specification (press release, December 2008)}. \url{https://www.khronos.org/news/press/the_khronos_group_releases_opencl_1.0_specification}. Vendor press release. The vendor-neutral standard with CUDA's model and different names (work-items, work-groups), six months after Apple's proposal.
\item \textbf{Wasson (The Tech Report) (2006).} \textit{Nvidia's GeForce 8800 graphics processor (launch review)}. \url{https://www.techreport.com/reviews/2006q4/geforce-8800/index.x?pg=1}. Press, secondary source. The launch-day architecture walkthrough with NVIDIA's briefing numbers, including the 185 W figure given to reviewers (155 W became the official maximum board power).
\end{itemize}
\subsection*{4. The CUDA model and its mapping to the chip (slides 9-10)}
\begin{itemize}[leftmargin=1.2em,itemsep=3pt]\sloppy
\item \textbf{Kirk and Hwu (2016).} \textit{Programming Massively Parallel Processors: A Hands-on Approach, 3rd edition}. \url{https://www.elsevier.com/books/programming-massively-parallel-processors/kirk/978-0-12-811986-0}. Textbook. The course that goes with slides 9-10 and 14: the CUDA model, memory hierarchy, tiling, and performance considerations; the 4th edition (Hwu, Kirk and El Hajj, 2022) adds later architectures.
\item \textbf{Harris (2007).} \textit{Optimizing parallel reduction in CUDA}. \url{https://developer.download.nvidia.com/compute/cuda/1.1-Beta/x86_website/projects/reduction/doc/reduction.pdf}. Vendor tutorial, open. Seven versions of one kernel, each removing a specific inefficiency (divergence, bank conflicts, idle threads, instruction overhead); the single best introduction to GPU performance thinking.
\item \textbf{Volkov and Demmel (2008).} \textit{Benchmarking GPUs to tune dense linear algebra}. \url{https://bebop.cs.berkeley.edu/pubs/volkov2008-benchmarking.pdf}. SC 2008 (Test of Time award, 2019), open copy. Microbenchmarks of G80 and GT200 used to design an SGEMM at 60\% of peak; the register-blocking lesson every later GEMM followed.
\item \textbf{Wittenbrink, Kilgariff and Prabhu (2011).} \textit{Fermi GF100 GPU architecture}. \url{https://doi.org/10.1109/MM.2011.24}. IEEE Micro 31(2). The Fermi SM of slide 10 described by its architects: dual warp schedulers, the configurable L1/shared array, the L2, ECC, and the graphics pipeline around them.
\item \textbf{NVIDIA (2009).} \textit{NVIDIA's next generation CUDA compute architecture: Fermi (whitepaper v1.1)}. \url{https://www.ece.lsu.edu/koppel/gp/refs/NVIDIA_Fermi_Compute_Architecture_Whitepaper.pdf}. Vendor whitepaper (mirror; the nvidia.com original is the same document). The G80/GT200/Fermi comparison table, the SM diagram, the memory hierarchy and the 48 warps per SM that slide 10 quotes.
\item \textbf{Glaskowsky (2009).} \textit{NVIDIA's Fermi: the first complete GPU computing architecture}. \url{https://www.nvidia.com/content/PDF/fermi_white_papers/P.Glaskowsky_NVIDIAFermi-TheFirstCompleteGPUComputingArchitecture.pdf}. Analyst whitepaper commissioned by NVIDIA; secondary source. Places Fermi against CPUs and earlier GPUs with unusual clarity about what 'general-purpose' required.
\end{itemize}
\subsection*{5. SIMT, latency hiding and the energy argument (slides 11-13)}
\begin{itemize}[leftmargin=1.2em,itemsep=3pt]\sloppy
\item \textbf{Fung, Sham, Yuan and Aamodt (2007).} \textit{Dynamic warp formation and scheduling for efficient GPU control flow}. \url{https://people.ece.ubc.ca/aamodt/papers/wwlfung.micro2007.pdf}. MICRO 2007, open copy. The first academic analysis of branch divergence in SIMT and a hardware proposal to regroup threads across warps; the cost model of slide 11.
\item \textbf{Hong and Kim (2009).} \textit{An analytical model for a GPU architecture with memory-level and thread-level parallelism awareness}. \url{https://doi.org/10.1145/1555754.1555775}. ISCA 2009. The MWP/CWP model: how many warps a memory system can serve at once and how many the compute pipeline needs; the formal version of slide 12.
\item \textbf{Bakhoda, Yuan, Fung, Wong and Aamodt (2009).} \textit{Analyzing CUDA workloads using a detailed GPU simulator}. \url{https://people.ece.ubc.ca/aamodt/papers/gpgpusim.ispass09.pdf}. ISPASS 2009, open copy. GPGPU-Sim, the simulator most GPU architecture research used for a decade, and its first findings on interconnect and memory bottlenecks.
\item \textbf{Rogers, O'Connor and Aamodt (2012).} \textit{Cache-conscious wavefront scheduling}. \url{https://engineering.purdue.edu/tgrogers/publication/rogers-micro-2012/rogers-micro-2012.pdf}. MICRO 2012, open copy. Why more resident warps can hurt: L1 thrashing, and a scheduler that limits the active set; the counterpoint to maximal occupancy.
\item \textbf{Volkov (2010).} \textit{Better performance at lower occupancy (GTC 2010 talk)}. \url{https://www.nvidia.com/content/GTC-2010/pdfs/2238_GTC2010.pdf}. Vendor conference slides, secondary source. Instruction-level parallelism within a thread as a substitute for thread-level parallelism; the argument slide 12 summarizes.
\item \textbf{Volkov (2016).} \textit{Understanding latency hiding on GPUs (PhD thesis, UC Berkeley EECS-2016-143)}. \url{https://www2.eecs.berkeley.edu/Pubs/TechRpts/2016/EECS-2016-143.pdf}. Open access. The rigorous version of Little's law for GPUs, measured on Fermi through Maxwell; the reference for anyone modeling occupancy.
\item \textbf{Keckler, Dally, Khailany, Garland and Glasco (2011).} \textit{GPUs and the future of parallel computing}. \url{https://www.cs.toronto.edu/~pekhimenko/courses/csc2231-f17/Papers/GPUs.pdf}. IEEE Micro 31(5), open copy (DOI 10.1109/MM.2011.89). The energy tables behind slide 13: 50 pJ per double-precision FMA at 40 nm, 14 pJ per 64-bit SRAM read, 240 fJ per bit-millimeter of wire, 28-45 pJ per bit of DRAM access and interface, and the Echelon design that followed from them.
\item \textbf{Esmaeilzadeh, Blem, St. Amant, Sankaralingam and Burger (2011).} \textit{Dark silicon and the end of multicore scaling}. \url{https://doi.org/10.1145/2000064.2000108}. ISCA 2011. At fixed power only part of a chip can be active; the argument for specialization that the tensor core later embodied.
\item \textbf{Dennard, Gaensslen, Yu, Rideout, Bassous and LeBlanc (1974).} \textit{Design of ion-implanted MOSFETs with very small physical dimensions}. \url{https://doi.org/10.1109/JSSC.1974.1050511}. IEEE JSSC 9(5). The scaling law whose breakdown around 2005 ended clock-speed growth and made throughput architectures inevitable.
\item \textbf{Bohr (2007).} \textit{A 30 year retrospective on Dennard's MOSFET scaling paper}. \url{https://doi.org/10.1109/N-SSC.2007.4785534}. IEEE SSCS Newsletter 12(1). Short, readable account of why the law stopped holding (leakage, threshold voltage) and what replaced it.
\item \textbf{Lee et al. (Intel) (2010).} \textit{Debunking the 100X GPU vs. CPU myth: an evaluation of throughput computing on CPU and GPU}. \url{https://www.cs.umd.edu/class/fall2019/cmsc714/readings/Lee-GPU.pdf}. ISCA 2010, open copy. Fourteen kernels tuned on both a GTX 280 and a Core i7-960: 2.5x on average. The paper that made honest comparisons the norm; slide 27.
\end{itemize}
\subsection*{6. The memory system and the roofline (slides 14-15)}
\begin{itemize}[leftmargin=1.2em,itemsep=3pt]\sloppy
\item \textbf{Williams, Waterman and Patterson (2009).} \textit{Roofline: an insightful visual performance model for multicore architectures}. \url{https://doi.org/10.1145/1498765.1498785}. CACM 52(4); open copy at escholarship.org/content/qt78h8v7mr/qt78h8v7mr.pdf. The model of slide 15: attainable performance as the minimum of peak and bandwidth times intensity.
\item \textbf{Wong, Papadopoulou, Sadooghi-Alvandi and Moshovos (2010).} \textit{Demystifying GPU microarchitecture through microbenchmarking}. \url{https://www.eecg.utoronto.ca/~moshovos/CUDA08/arx/microbenchmark_report.pdf}. ISPASS 2010 (DOI 10.1109/ISPASS.2010.5452013), open technical-report copy. Pointer-chasing and dependent-chain measurements on GT200: about 440 cycles to global memory, 38 to shared memory, a 24-cycle arithmetic pipeline, the cache and TLB structure.
\item \textbf{Mei and Chu (2017).} \textit{Dissecting GPU memory hierarchy through microbenchmarking}. arXiv:1509.02308. IEEE TPDS 28(1). The same method applied to Fermi, Kepler and Maxwell: cache sizes, set structure, latencies; fills the gap between Wong (GT200) and Jia (Volta).
\item \textbf{JEDEC (2016).} \textit{JEDEC updates High Bandwidth Memory standard JESD235A (press release, January 12, 2016)}. \url{https://www.jedec.org/node/5916}. Standards body press release. HBM2: eight independent 128-bit channels per stack, up to 256 GB/s and 8 GB per stack, pseudo-channel mode; the memory of Pascal and Volta.
\item \textbf{JEDEC (2016).} \textit{JEDEC announces publication of GDDR5X graphics memory standard JESD232 (January 21, 2016)}. \url{https://www.jedec.org/node/5917}. Standards body press release. The other 2016 path to bandwidth, used on the consumer GP104 instead of HBM2.
\item \textbf{NVIDIA (2016).} \textit{GeForce GTX 1080 whitepaper (GP104)}. \url{https://kolegite.com/EE_library/datasheets_and_manuals/GPU/GeForce_GTX_1080_Whitepaper_FINAL.pdf}. Vendor whitepaper (mirror). The consumer Pascal with GDDR5X; its comparison table is also the accessible source for GM204 (GTX 980) figures.
\end{itemize}
\subsection*{7. The generations: Tesla, Fermi, Kepler, Maxwell, Pascal (slides 16-21)}
\begin{itemize}[leftmargin=1.2em,itemsep=3pt]\sloppy
\item \textbf{NVIDIA (2010).} \textit{Tesla C1060 computing processor board specification (BD-04111-001\_v06)}. \url{https://www.nvidia.com/docs/io/43395/bd-04111-001_v06.pdf}. Vendor datasheet. GT200 as a compute product: 240 cores at 1.296 GHz, 4 GB GDDR3 at 102 GB/s, 187.8 W.
\item \textbf{NVIDIA (2010).} \textit{Tesla C2050 / C2070 GPU computing processor datasheet (July 2010)}. \url{https://www.nvidia.com/docs/io/43395/nv_ds_tesla_c2050_c2070_jul10_lores.pdf}. Vendor datasheet. Fermi as shipped: 448 cores, 515 GFLOP/s double, 1.03 TFLOP/s single, 144 GB/s, 238 W.
\item \textbf{NVIDIA (2012).} \textit{NVIDIA's next generation CUDA compute architecture: Kepler GK110 (whitepaper v1.0)}. \url{https://www.nvidia.co.jp/content/apac/pdf/tesla/nvidia-kepler-gk110-architecture-whitepaper-jp.pdf}. Vendor whitepaper (mirror of the English document). The SMX, dynamic parallelism, Hyper-Q, GPUDirect RDMA, and the compute-capability table from Fermi to GK110.
\item \textbf{NVIDIA (2014).} \textit{Tesla K80 / K40 datasheet}. \url{https://hpc-test.llnl.gov/sites/default/files/Tesla-K80K40-datasheet.pdf}. Vendor datasheet (mirror). K40 at 4.29 / 1.43 TFLOP/s base and 5.04 / 1.68 boost; K80 at 8.74 / 2.91 with boost, 24 GB, 480 GB/s.
\item \textbf{NVIDIA (2016).} \textit{Tesla M40 datasheet (January 2016)}. \url{https://images.nvidia.com/content/tesla/pdf/nvidia-teslam40-datasheet.pdf}. Vendor datasheet. Maxwell for training: 3,072 cores, 7 TFLOP/s with boost, 288 GB/s, 250 W, and the FP64 rate that kept it out of HPC.
\item \textbf{NVIDIA (2016).} \textit{NVIDIA Tesla P100 whitepaper (WP-08019-001\_v01.1)}. \url{https://techpowerup.com/gpu-specs/docs/nvidia-gp100-architecture.pdf}. Vendor whitepaper (mirror of the same document). GP100 in full: the SM, FP16, HBM2 and CoWoS, NVLink 1.0, unified memory with page faulting, compute preemption, and Table 1 comparing K40, M40 and P100. Slide 20.
\item \textbf{NVIDIA (2016).} \textit{NVIDIA delivers massive performance leap for deep learning, HPC applications with NVIDIA Tesla P100 accelerators (press release, April 5, 2016)}. \url{https://www.globenewswire.com/news-release/2016/04/05/1149699/0/en/NVIDIA-Delivers-Massive-Performance-Leap-for-Deep-Learning-HPC-Applications-With-NVIDIA-Tesla-P100-Accelerators.html}. Vendor press release. The launch claims as made: 5.3 / 10.6 / 21.2 TFLOP/s, 720 GB/s, 160 GB/s NVLink, 15.3 billion transistors.
\item \textbf{Foley and Danskin (2017).} \textit{Ultra-performance Pascal GPU and NVLink interconnect}. \url{https://doi.org/10.1109/MM.2017.37}. IEEE Micro 37(2). GP100 and NVLink by their architects; the Hot Chips 28 talk by the same authors (old.hotchips.org, HC28.22.121) has the floorplan and PHY details.
\item \textbf{NVIDIA (2015).} \textit{New NVIDIA hyperscale accelerators boost machine learning throughput for web data centers (press release, November 10, 2015)}. \url{https://www.globenewswire.com/news-release/2015/11/10/1149678/0/en/New-NVIDIA-Hyperscale-Accelerators-Boost-Machine-Learning-Throughput-for-Web-Data-Centers.html}. Vendor press release. The M40 announcement, and the moment NVIDIA named deep-learning training as a product category of its own.
\item \textbf{Krizhevsky, Sutskever and Hinton (2012).} \textit{ImageNet classification with deep convolutional neural networks}. \url{https://papers.nips.cc/paper/4824-imagenet-classification-with-deep-convolutional-neural-networks.pdf}. NIPS 2012, open access (also CACM 60(6), 2017). Trained 'between five and six days on two GTX 580 3GB GPUs'; the paper that redirected the architecture. Slide 19.
\item \textbf{Chetlur, Woolley, Vandermersch, Cohen, Tran, Catanzaro and Shelhamer (2014).} \textit{cuDNN: efficient primitives for deep learning}. arXiv:1410.0759. The library that tied every deep-learning framework to CUDA; convolution as GEMM with tiling in on-chip memory, the kernels the tensor core later accelerated.
\item \textbf{Raina, Madhavan and Ng (2009).} \textit{Large-scale deep unsupervised learning using graphics processors}. \url{https://icml.cc/imls/conferences/2009/papers/218.pdf}. ICML 2009, open copy. Deep belief networks and sparse coding on a GTX 280: the first widely read case that GPUs change what deep learning can afford.
\item \textbf{Ciresan, Meier, Gambardella and Schmidhuber (2010).} \textit{Deep, big, simple neural nets for handwritten digit recognition}. arXiv:1003.0358. Neural Computation 22(12). Plain multilayer perceptrons trained on GPUs beat the MNIST state of the art; evidence that compute, not cleverness, was the missing ingredient.
\item \textbf{Tom's Hardware (Angelini) (2015).} \textit{Nvidia GeForce GTX Titan X 12GB review (March 17, 2015)}. \url{https://www.tomshardware.com/reviews/nvidia-geforce-gtx-titan-x-gm200-maxwell,4091.html}. Press, secondary source. The GM200 launch: 8 billion transistors, 601 mm2, 3,072 cores, FP64 at 1/32; the consumer die that became the Tesla M40.
\end{itemize}
\subsection*{8. Volta and the tensor core (slides 22-24)}
\begin{itemize}[leftmargin=1.2em,itemsep=3pt]\sloppy
\item \textbf{NVIDIA (2017).} \textit{NVIDIA Tesla V100 GPU architecture (whitepaper WP-08608-001\_v1.1)}. \url{https://www.ece.lsu.edu/gp/refs/volta-architecture-whitepaper.pdf}. Vendor whitepaper (mirror of the same document). The GV100 SM with its four processing blocks, tensor cores, independent thread scheduling, the unified L1/shared memory, NVLink 2.0, and Table 1 comparing P100 and V100. Slides 21-24.
\item \textbf{Choquette, Giroux and Foley (2018).} \textit{Volta: performance and programmability}. \url{https://doi.org/10.1109/MM.2018.022071134}. IEEE Micro 38(2). Volta by its architects: why the SM was repartitioned, how tensor cores fit the pipeline, and what independent thread scheduling cost and bought.
\item \textbf{NVIDIA (2018).} \textit{Tesla V100 data sheet (March 2018)}. \url{https://images.nvidia.com/content/technologies/volta/pdf/tesla-volta-v100-datasheet-letter-fnl-web.pdf}. Vendor datasheet. SXM2 against PCIe: 7.8 / 15.7 / 125 against 7 / 14 / 112 TFLOP/s, 300 against 250 W, 16 or 32 GB.
\item \textbf{Jia, Maggioni, Staiger and Scarpazza (2018).} \textit{Dissecting the NVIDIA Volta GPU architecture via microbenchmarking}. arXiv:1804.06826. Citadel's measurement of V100: 19-cycle shared memory, 28-cycle L1, about 193-cycle L2, 375-cycle global memory with a TLB hit and 1,029 with a miss, 4-cycle FP32 FMA, tensor-core behavior, and the instruction encoding. Slides 12 and 24.
\item \textbf{Markidis, Der Chien, Laure, Peng and Vetter (2018).} \textit{NVIDIA tensor core programmability, performance and precision}. arXiv:1803.04014. IPDPSW 2018. The WMMA API measured: achieved TFLOP/s against the 125 peak, and the precision lost to FP16 inputs with and without FP32 accumulation.
\item \textbf{Micikevicius et al. (2017).} \textit{Mixed precision training}. arXiv:1710.03740. ICLR 2018. FP32 master weights, FP16 forward and backward passes, loss scaling; the recipe that made tensor cores usable for training without accuracy loss. Slide 23.
\item \textbf{NVIDIA (via Edge AI and Vision Alliance) (2017).} \textit{NVIDIA launches revolutionary Volta GPU platform (press release, May 10, 2017)}. \url{https://www.edge-ai-vision.com/2017/05/nvidia-launches-revolutionary-volta-gpu-platform-fueling-next-era-of-ai-and-high-performance-computing/}. Vendor press release (mirror). The launch claims: 21 billion transistors, 640 tensor cores, 120 TFLOP/s (later 125), 900 GB/s.
\end{itemize}
\subsection*{9. Systems: interconnect, supercomputers, the alternatives (slides 25-27)}
\begin{itemize}[leftmargin=1.2em,itemsep=3pt]\sloppy
\item \textbf{Danskin and Foley (2016).} \textit{Pascal GPU with NVLink (Hot Chips 28 slides)}. \url{https://old.hotchips.org/wp-content/uploads/hc_archives/hc28/HC28.22-Monday-Epub/HC28.22.10-GPU-HPC-Epub/HC28.22.121-Pascal-GPU-DanskinFoley-NVIDIA-v06-6_7.pdf}. Conference slides, open. The GP100 floorplan, the HBM2 and CoWoS package, and the NVLink physical layer and the DGX-1 topology. Slide 25.
\item \textbf{Ishii and Wells (NVIDIA) (2018).} \textit{NVSwitch and DGX-2 (Hot Chips 30 slides)}. \url{https://www.old.hotchips.org/hc30/2conf/2.01_Nvidia_NVswitch_HotChips2018_DGX2NVS_Final.pdf}. Conference slides, open. The 18-port NVSwitch (2 billion transistors, 106 mm2, 900 GB/s aggregate) and the DGX-2's 16-GPU fabric with 2.4 TB/s of bisection bandwidth.
\item \textbf{HPCwire (2018).} \textit{Nvidia ups hardware game with 16-GPU DGX-2 server and NVSwitch (March 27, 2018)}. \url{https://www.hpcwire.com/2018/03/27/nvidia-ups-hardware-game-with-16-gpu-dgx-2-server-and-nvswitch/}. Press, secondary source. The DGX-2 and V100 32 GB announcement with the figures as given at GTC 2018.
\item \textbf{Oak Ridge Leadership Computing Facility (2018).} \textit{Summit by the numbers (June 2018)}. \url{https://www.olcf.ornl.gov/wp-content/uploads/2018/06/Summit_bythenumbers_FIN-1.pdf}. Facility fact sheet. 4,608 nodes, 27,648 V100s, 200 PFLOP/s peak, 13 MW, the node layout of slide 26.
\item \textbf{Lawrence Livermore National Laboratory (2018).} \textit{Sierra platform page}. \url{https://hpc.llnl.gov/hardware/platforms/sierra}. Facility page. 4,320 nodes, 17,280 V100s, 125.6 PFLOP/s peak: Summit's sibling with four GPUs per node.
\item \textbf{TOP500 (2010).} \textit{TOP500 list, November 2010}. \url{https://top500.org/lists/top500/2010/11/}. Benchmark list. Tianhe-1A at 2.566 PFLOP/s, the first GPU-accelerated number one; TSUBAME 2.0 and Nebulae in the top five.
\item \textbf{NVIDIA (2010).} \textit{Three of the world's top five supercomputers powered by NVIDIA Tesla GPUs (press release, November 15, 2010)}. \url{https://www.globenewswire.com/news-release/2010/11/15/1149888/0/en/Three-of-the-World-s-Top-Five-Supercomputers-Powered-by-NVIDIA-Tesla-GPUs.html}. Vendor press release. 7,168 GPUs in Tianhe-1A and the Fermi-era systems; note that it quotes the pre-list 2.507 PFLOP/s figure.
\item \textbf{TOP500 (2012).} \textit{TOP500 list, November 2012}. \url{https://top500.org/lists/top500/2012/11/}. Benchmark list. Titan at 17.59 PFLOP/s with 18,688 K20Xs, number one.
\item \textbf{TOP500 (2017).} \textit{TOP500 list, November 2017}. \url{https://top500.org/lists/top500/list/2017/11}. Benchmark list. Piz Daint at 19.59 PFLOP/s with 5,320 P100s, number three; Titan still number five.
\item \textbf{TOP500 (2018).} \textit{TOP500 list, June 2018}. \url{https://www.top500.org/lists/2018/06/}. Benchmark list. Summit's first appearance at 122.3 PFLOP/s, number one; Sierra at 71.6.
\item \textbf{TOP500 (2018).} \textit{TOP500 list, November 2018}. \url{https://top500.org/lists/2018/11/}. Benchmark list. Summit at 143.5 PFLOP/s, Sierra at 94.64, Piz Daint at 21.23; the 148.6 figure often quoted for Summit is from June 2019.
\item \textbf{TOP500 News (2017).} \textit{Swiss deploy world's fastest GPU-powered supercomputer}. \url{https://www.top500.org/news/swiss-deploy-worlds-fastest-gpu-powered-supercomputer}. Benchmark-site news, secondary source. The Piz Daint P100 upgrade's timeline and node count.
\item \textbf{Seiler et al. (Intel) (2008).} \textit{Larrabee: a many-core x86 architecture for visual computing}. \url{https://doi.org/10.1145/1360612.1360617}. ACM TOG 27(3), SIGGRAPH 2008; open copy at courses.cs.washington.edu/courses/cse591n/08au/Papers/larrabee\_manycore.pdf. The x86 alternative: many small cores, 16-wide vectors, coherent caches, graphics in software. Slide 27.
\item \textbf{Feldman (HPCwire) (2009).} \textit{Intel cancels 2010 Larrabee debut (December 7, 2009)}. \url{https://www.hpcwire.com/2009/12/07/intel_cancels_2010_larrabee_debut/}. Press, secondary source. Intel's December 4, 2009 statement that Larrabee would not ship as a discrete graphics product.
\item \textbf{Feldman (HPCwire) (2012).} \textit{Intel brings manycore x86 to market with Knights Corner (November 12, 2012)}. \url{https://www.hpcwire.com/2012/11/12/intel_brings_manycore_x86_to_market_with_knights_corner/}. Press, secondary source. The Xeon Phi 5110P: 60 cores, 1.011 TFLOP/s double, 8 GB, 320 GB/s, 225 W (the article misprints the core count).
\item \textbf{Feldman (TOP500 News) (2016).} \textit{Intel takes on NVIDIA with Knights Landing launch (June 2016)}. \url{https://www.top500.org/news/intel-takes-on-nvidia-with-knights-landing-launch/}. Secondary source. Knights Landing SKUs (up to 72 cores, about 3.46 TFLOP/s double) and the 16 GB of MCDRAM at 490 GB/s.
\item \textbf{HPCwire (2018).} \textit{End of the road for Knights Landing Phi (July 25, 2018)}. \url{https://www.hpcwire.com/2018/07/25/end-of-the-road-for-knights-landing-phi}. Press, secondary source. The discontinuation notice: last orders August 2018, last shipments July 2019.
\item \textbf{AMD (2014).} \textit{AMD unleashes world's most powerful server GPU for HPC (press release, August 6, 2014)}. \url{https://www.globenewswire.com/news-release/2014/08/06/1011324/0/en/AMD-Unleashes-World-s-Most-Powerful-Server-GPU-for-HPC.html}. Vendor press release. The FirePro S9150: 2.53 TFLOP/s double, 5.07 single, 16 GB, 320 GB/s, 235 W; GCN hardware ahead of NVIDIA's on paper.
\item \textbf{AMD (2015).} \textit{AMD launches 'Boltzmann Initiative' to dramatically reduce barriers to GPU computing on AMD FirePro graphics (press release, November 16, 2015)}. \url{https://ir.amd.com/news-events/press-releases/detail/652/amd-launches-boltzmann-initiative-to-dramatically-reduce-barriers-to-gpu-computing-on-amd-fireprotm-graphics}. Vendor press release. HCC, HIP and the headless Linux driver that became ROCm; the start of AMD's CUDA-compatible software path.
\item \textbf{AMD (2016).} \textit{AMD releases new version of ROCm, the most versatile open source platform for GPU computing (press release, November 14, 2016)}. \url{https://ir.amd.com/news-events/press-releases/detail/732/amd-releases-new-version-of-rocm-the-most-versatile-open-source-platform-for-gpu-computing}. Vendor press release. ROCm 1.3 at SC16; ROCm 1.0 had shipped in April 2016.
\end{itemize}
\subsection*{10. Practice, textbooks and retrospectives (slides 28-31)}
\begin{itemize}[leftmargin=1.2em,itemsep=3pt]\sloppy
\item \textbf{Halfhill (Microprocessor Report) (2009).} \textit{Looking beyond graphics (Fermi analysis)}. \url{https://www.halfhill.com/mpr/2009/234001.pdf}. Analyst report, open copy; secondary source. A contemporaneous independent reading of Fermi's design choices and their cost.
\item \textbf{NVIDIA (2014).} \textit{Kepler GK110/210 compute architecture whitepaper}. \url{https://techpowerup.com/gpu-specs/docs/nvidia-gk110-210-compute-architecture.pdf}. Vendor whitepaper (mirror). The GK210 revision for the K80: doubled register file and shared memory per SMX.
\item \textbf{NVIDIA (2017).} \textit{CUDA C programming guide (version 9.0)}. \url{https://docs.nvidia.com/cuda/archive/9.0/cuda-c-programming-guide/index.html}. Vendor documentation. The compute-capability tables (registers, shared memory, warps per SM for every generation), the memory model, and the Volta-era additions (cooperative groups, WMMA, independent thread scheduling).
\item \textbf{Hwu, Kirk and El Hajj (2022).} \textit{Programming Massively Parallel Processors: A Hands-on Approach, 4th edition}. \url{https://www.elsevier.com/books/programming-massively-parallel-processors/hwu/978-0-323-91231-0}. Textbook. The current edition of the CUDA course, with Volta and later architectures; the performance chapters are the practice column of slide 29 at book length.
\item \textbf{Hennessy and Patterson (2011).} \textit{Computer Architecture: A Quantitative Approach, 5th edition, chapter 4 and appendix}. \url{https://www.elsevier.com/books/computer-architecture/hennessy/978-0-12-383872-8}. Textbook. The edition whose GPU section is built on Fermi, with the vector-processor vocabulary that makes SIMT precise.
\item \textbf{Peddie (IEEE Computer Society, Chasing Pixels) (2020).} \textit{Famous graphics chips: Geometry Engine}. \url{https://www.computer.org/publications/tech-news/chasing-pixels/geometry-engine}. Secondary source. A short illustrated history of the Geometry Engine for readers who want the 1982 story in plain language.
\end{itemize}
